\sct{Series and Improper Integral Theorems and Methods}
\ssc{Limit of real series theorems}
\sssc{Integral test (for convergence) or Maclaurin–Cauchy test}
Suppose $f(x)\colon[a,\infty)\to\bbR$ for some $a\in\bbN$ is monotone and $\lim_{x\to\infty}f(x)=0$. Then $\int_a^{\infty}f(x)\dd{x}$ is convergent iff $\sum_{k=a}^{\infty}f(k)$ is convergent. And if $\int_a^{\infty}f(x)\dd{x}$ is convergent and $f(x)$ is monotone decreasing, for any $n\in\bbN_{\ge a}$,
\[-f(n)+\int_n^{\infty}f(x)\dd{x}\le\sum_{k=n+1}^{\infty}f(k)\le\int_n^{\infty}f(x)\dd{x}\le\sum_{k=n}^{\infty}f(k)\le f(n)+\int_n^{\infty}f(x)\dd{x},\]
called remainder estimate for integral test;

if $\int_a^{\infty}f(x)\dd{x}$ is convergent and $f(x)$ is monotone increasing, for any $n\in\bbN_{\ge a}$,
\[-f(n)+\int_n^{\infty}f(x)\dd{x}\ge\sum_{k=n+1}^{\infty}f(k)\ge\int_n^{\infty}f(x)\dd{x}\ge\sum_{k=n}^{\infty}f(k)\ge f(n)+\int_n^{\infty}f(x)\dd{x},\]
called remainder estimate for integral test.
\bpr
Consider the case where $f(x)$ is monotone increasing. The case where $f(x)$ is monotone decreasing can be proven similarly.

First consider the case where $\int_a^{\infty}f(x)\dd{x}$ is convergent. This implies that $f(x)$ is non-positive. Thus, for any fixed $n\in\bbN_{\ge a},
\[\int_n^{n+1}f(x)\dd{x}\ge f(n),\]
and for any $k\geq n+1$,
\[\int_k^{k+1}f(x)\dd{x}\ge f(k)\ge\int_{k-1}^kf(x)\dd{x}.\]
For any $N\in\bbN_{\ge n}$, summing over all $k\in\bbN_{\ge n}$, we obtain, 
\[\int_n^{N+1}f(x)\dd{x}\ge\sum_{k=n}^Nf(k)\ge f(n)+\int_n^Nf(x)\dd{x}.\]
Take limit $\lim_{N\to\infty}$, we obtain
\[\int_n^{\infty}f(x)\dd{x}\ge\sum_{k=n}^{\infty}f(n)\ge f(n)+\int_n^{\infty}f(x)\dd{x}.\]
Second, consider the case where $\int_a^{\infty}f(x)\dd{x}$ is divergent. Since $f(x)$ is monotone. Either $\int_a^{\infty}f(x)\dd{x}=-\infty$ or $\int_a^{\infty}f(x)\dd{x}=\infty$. For the former, the supremum of $f(x)$ is negative, and $\sum_{k=a}^{\infty}f(k)=-\infty$. For the latter, there exists $x_0>a$ such that $f(x)>0$ for all $x\ge x_0$, and thus $\sum_{k=a}^{\infty}f(k)=\infty$.
\epr
\sssc{Alternating series test or Leibniz's test}
If a series $\sum_{n=1}^{\infty}(-1)^{n-1}b_n$ satisfies that
\[b_n\geq b_{n+1}\geq 0,\quad n\in\bbN,\]
\[\lim_{n\to\infty}b_n=0,\]
then it is convergent, and
\[0\leq\abs{\sum_{k=n}^{\infty}(-1)^{k-1}b_k}\leq b_n,\forall n\in\bbN,\]
where the latter is called alternating series estimation theorem.
\bpr
Let
\[c_1=b_1,\]
\[c_n=b_{2n-1}-b_{2n-2},\quad n\in\bbN\land n>1,\]
\[S_n=\sum_{k=1}^n(-1)^{k-1}b_k,\]
\[S=\sum_{n=1}^{\infty}(-1)^{n-1}b_n,\]
\[S_{2n-1}=T_n=\sum_{k=1}^nc_n.\]
$T_n$ is non-increasing since
\[c_n\leq 0,\quad\forall n\in\bbN\land n>1.\]
$T_n$ is bounced-below by $0$ since
\[c_1\geq b_2\geq-\sum_{k=2}^{\infty}c_k,\]
\[T_n\geq 0.\]
By monotone convergence theorem, $T_n$ is convergent, and thus $S_n$ is convergent.
\[b_{2n-2}\geq-\sum_{k=n}^{\infty}c_k=S-T_{n-1}=S-S_{2n-3}.\]
Thus for even $n$,
\[b_n\geq\sum_{k=n}^{\infty}(-1)^{k-1}b_k=\abs{\sum_{k=n}^{\infty}(-1)^{k-1}b_k}.\]
For odd $n$,
\[\abs{\sum_{k=n}^{\infty}(-1)^{k-1}b_k}=-\sum_{k=n}^{\infty}(-1)^{k-1}b_k=-b_n+\sum_{k=n+1}^{\infty}(-1)^{k-1}b_k\geq -b_n.\]
\epr
\sssc{Ratio test}
\ben
\item If $\lim_{n\to\infty}\abs{\frac{a_{n+1}}{a_n}}=L<1$, then the series $\sum_{n=1}^{\infty}a_n$ is absolutely convergent.
\item If $\lim_{n\to\infty}\abs{\frac{a_{n+1}}{a_n}}=L>1$ or $\lim_{n\to\infty}\abs{\frac{a_{n+1}}{a_n}}=\infty$, then the series $\sum_{n=1}^{\infty}a_n$ is divergent.
\een
\bpr
For the first statement, choose an arbitrary $L<r<1$. There exists an integer $N$ such that
\[\abs{\frac{a_{n+1}}{a_n}}<r,\quad\forall n\geq N.\]
Thus,
\[\abs{a_{N+k}}<\abs{a_N}r^k,\quad\forall k\geq 1.\]
Since
\[\sum_{k=1}^{\infty}\abs{a_N}r^k\]
is convergent, by comparison test, $\sum_{n=1}^{\infty}\abs{a_n}$ is convergent.

For the second statement,
\[\lim_{n\to\infty}a_n\neq 0,\]
thus $\sum_{n=1}^{\infty}a_n$ is divergent.
\epr
\sssc{(Cauchy) root test}
\ben
\item If $\limsup_{n\to\infty}\sqrt[n]{\abs{a_n}}<1$, then the series $\sum_{n=1}^{\infty}a_n$ is absolutely convergent.
\item If $\limsup_{n\to\infty}\sqrt[n]{\abs{a_n}}>1$, then the series $\sum_{n=1}^{\infty}a_n$ is divergent.
\een
\bpr
Let
\[L=\limsup_{n\to\infty}\sqrt[n]{\abs{a_n}}.\]
For the first statement, choose an arbitrary $L<r<1$. There exists an integer $N$ such that
\[\sqrt[n]{\abs{a_n}}<r,\quad\forall n\geq N.\]
Thus,
\[\abs{a_n}<r^n,\quad\forall n\geq N.\]
Since
\[\sum_{n=1}^{\infty}r^n\]
is convergent, by comparison test, $\sum_{n=1}^{\infty}\abs{a_n}$ is convergent.

For the second statement,
\[\neg\lim_{n\to\infty}a_n=0.\]
Thus $\sum_{n=1}^{\infty}a_n$ is divergent.
\epr
\sssc{Riemann series theorem or Riemann rearrangement theorem}
Let $\langle a_n\rangle_{n=1}^{\infty}$ be a real sequence, and that $\sum_{n=1}^{\infty}a_n$ is conditionally convergent. Then, for any extended-real number $M$, there exists a permutation $\sigma$ on $\bbN$ such that
\[\sum_{n=1}^{\infty}a_{\sigma(n)}=M.\]
\bpr
Since $\sum_{n=1}^{\infty}a_n$ is conditionally convergent, $\langle a_n\rangle$ has infinitely many positive terms and infinitely many negative terms. Let
\[a_n^+=\bcs a_n,\quad&a_n>0\\0,\quad&a_n\leq 0\ecs,\]
\[a_n^-=\bcs a_n,\quad&a_n<0\\0,\quad&a_n\geq 0\ecs.\]
Since $\sum_{n=1}^{\infty}a_n$ is conditionally convergent,
\[\sum_{n=1}^{\infty}a_n^+=\infty,\]
\[\sum_{n=1}^{\infty}a_n^-=-\infty.\]
For any real number $M$, let $p_1$ be the smallest positive integer such that
\[M<\sum_{n=1}^{p_1}a_n^+,\]
and let $q_1$ be the smallest positive integer such that
\[M>\sum_{n=1}^{p_1}a_n^++\sum_{n=1}^{q_1}a_n^-.\]
Now recursively, for all $k\in\bbN_{>1}$, let $p_k$ be the smallest integer such that $p_k>p_{k-1}$ and
\[M<\sum_{n=1}^{p_k}a_n^++\sum_{n=1}^{q_{k-1}}a_n^-,\]
and let $q_k$ be the smallest integer such that $q_k>q_{k-1}$ and
\[M>\sum_{n=1}^{p_k}a_n^++\sum_{n=1}^{q_k}a_n^-.\]
The new sequence is defined as
\[\langle a_n^+\rangle_{n=1\land a_n\geq 0}^{p_1}\langle a_n^-\rangle_{n=1\land a_n<0}^{q_1}\langle\langle a_n^+\rangle_{n=p_k+1\land a_n\geq 0}^{p_{k+1}}\langle a_n^-\rangle_{n=q_k+1\land a_n<0}^{q_{k+1}}\rangle_{k=1}^{\infty}.\]
We have
\[\sum_{n=1}^{p_{k+1}-1}a_n^++\sum_{n=1}^{q_k}a_n^-\leq M<\sum_{n=1}^{p_{k+1}}a_n^++\sum_{n=1}^{q_k}a_n^-,\]
\[0<\qty(\sum_{n=1}^{p_{k+1}}a_n^++\sum_{n=1}^{q_k}a_n^-)-M\leq a_{p_{k+1}}^+.\]
Since $\sum_{n=1}^{\infty}a_n$ is conditionally convergent,
\[\lim_{k\to\infty}a_{p_{k+1}}^+=0.\]
By squeeze theorem,
\[\lim_{k\to\infty}\sum_{n=1}^{p_{k+1}}a_n^++\sum_{n=1}^{q_k}a_n^-=M.\]
Similarly,
\[\sum_{n=1}^{p_k}a_n^++\sum_{n=1}^{q_k-1}a_n^-\geq M>\sum_{n=1}^{p_k}a_n^++\sum_{n=1}^{q_k}a_n^-,\]
\[0\geq M-\qty(\sum_{n=1}^{p_k}a_n^++\sum_{n=1}^{q_k-1}a_n^-)>a_{q_k}^-.\]
Since $\sum_{n=1}^{\infty}a_n$ is convergent,
\[\lim_{k\to\infty}a_{q_k}^-=0.\]
By squeeze theorem,
\[\lim_{k\to\infty}\sum_{n=1}^{p_k}a_n^++\sum_{n=1}^{q_k}a_n^-=M.\]
Since for all $p_k<p<p_{k+1}$,
\[\sum_{n=1}^{p_k}a_n^++\sum_{n=1}^{q_k}a_n^-\leq\sum_{n=1}^pa_n^++\sum_{n=1}^{q_k}a_n^-\leq\sum_{n=1}^{p_{k+1}}a_n^++\sum_{n=1}^{q_k}a_n^-,\]
and for all $q_k<q<q_{k+1}$,
\[\sum_{n=1}^{p_{k+1}}a_n^++\sum_{n=1}^{q_k}a_n^-\geq\sum_{n=1}^{p_{k+1}}a_n^++\sum_{n=1}^qa_n^-\geq\sum_{n=1}^{p_{k+1}}a_n^++\sum_{n=1}^{q_{k+1}}a_n^-,\]
by squeeze theorem, the new sequence converges to $M$.

For $\infty$ case ($-\infty$ case can be proven similarly), let $p_1$ be the smallest positive integer such that
\[1<\sum_{n=1}^{p_1}a_n^+,\]
and let $q_1$ be the smallest positive integer such that
\[1>\sum_{n=1}^{p_1}a_n^++\sum_{n=1}^{q_1}a_n^-.\]
Now recursively, for all $k\in\bbN_{>1}$, let $p_k$ be the smallest integer such that $p_k>p_{k-1}$ and
\[k<\sum_{n=1}^{p_k}a_n^++\sum_{n=1}^{q_{k-1}}a_n^-,\]
and let $q_k$ be the smallest integer such that $q_k>q_{k-1}$ and
\[k>\sum_{n=1}^{p_k}a_n^++\sum_{n=1}^{q_k}a_n^-.\]
The new sequence is defined as
\[\langle a_n^+\rangle_{n=1\land a_n\geq 0}^{p_1}\langle a_n^-\rangle_{n=1\land a_n<0}^{q_1}\langle\langle a_n^+\rangle_{n=p_k+1\land a_n\geq 0}^{p_{k+1}}\langle a_n^-\rangle_{n=q_k+1\land a_n<0}^{q_{k+1}}\rangle_{k=1}^{\infty}.\]
Thus for arbitrary large $k$, we have
\[k<\sum_{n=1}^{p_k}a_n^++\sum_{n=1}^{q_{k-1}}a_n^-,\]
and
\[k<\sum_{n=1}^{p_k}a_n^-+\sum_{n=1}^{q_k-1}a_n^+.\]
Since for all $p_k<p<p_{k+1}$,
\[\sum_{n=1}^{p_k}a_n^++\sum_{n=1}^{q_k}a_n^-\leq\sum_{n=1}^pa_n^++\sum_{n=1}^{q_k}a_n^-\leq\sum_{n=1}^{p_{k+1}}a_n^++\sum_{n=1}^{q_k}a_n^-,\]
and for all $q_k-1<q<q_{k+1}-1$,
\[\sum_{n=1}^{p_{k+1}}a_n^++\sum_{n=1}^{q_k-1}a_n^-\geq\sum_{n=1}^{p_{k+1}}a_n^++\sum_{n=1}^qa_n^-\geq\sum_{n=1}^{p_{k+1}}a_n^++\sum_{n=1}^{q_{k+1}-1}a_n^-,\]
by squeeze theorem, the new sequence has limit $\infty$.
\epr
\sssc{Cauchy–Hadamard theorem}
For any given power series $\sum_{n=0}^{\infty}c_n(x-a)^n$, let
\[L=\limsup_{n\to\infty}\sqrt[n]{\abs{c_n}}\in\bar{\bbR}_{\geq 0},\]
then,
\[R=\bcs\infty,\quad&L=0\\\frac{1}{L},\quad&L\in\bbR_{>0}\\0,\quad&L=\infty\ecs,\]
called radius of convergence, is such that the series is absolute convergent if $|x-a|<R$ and divergent if $|x-a|>R$, and if $R=0$, the series is absolutely convergent at $x=a$ and divergent elsewhere.
\bpr
Let $r=|x-a|$.
\[\limsup_{n\to\infty}\sqrt[n]{\abs{c_n(x-a)^n}}=\bcs Lr,\quad&r\neq 0\\0,\quad&r=0\ecs.\]

If $L=0$, then the series is absolutely convergent for all $x$, thus $R=\infty$.

Else if $L\in\bbR_{>0}$, then the series is absolutely convergent for $r<\frac{1}{L}$ and divergent for $r>\frac{1}{L}$, thus $R=\frac{1}{L}$.

Else if $L=\infty$, then the series is absolutely convergent at $x=a$ and divergent elsewhere, thus $R=0$.
\epr
\sssc{Power series term-by-term differentiation and integration theorem}
If the power series $\sum_{n=0}^{\infty}c_n(x-a)^n$ has radius of convergence $R>0$, then the function $f$ defined as
\[f(x)=\sum_{n=0}^{\infty}c_n(x-a)^n\]
is differentiable on $(a-R,a+R)$, and
\[f'(x)=\sum_{n=0}^{\infty}nc_n(x-a)^{n-1},\]
\[\int f(x)\dd{x}=C+\sum_{n=0}^{\infty}c_n\frac{(x-a)^{n+1}}{n+1},\]
and the radii of convergence of both of them are $R$.
\bpr
PLACEHOLDER
\epr
\sssc{Cauchy's convergence test}
A series $\sum_{n=0}^{\infty}a_n$ converges iff for every $\varepsilon>0$, there is a natural number $N$ such that
\[\abs{\sum_{i=n+1}^{n+p}a_i}<\varepsilon\]
for all $n\geq N$ and $p\geq 1$.
\sssc{Mertens' theorem of Cauchy product}
If $\sum_{n=0}^{\infty}a_n$ converges to $A$, $\sum_{n=0}^{\infty}b_n$ converges to $B$, and at least one of them converges absolutely, then their Cauchy product converges to $AB$.

If $\sum_{n=0}^{\infty}a_n$ converges to $A$, $\sum_{n=0}^{\infty}b_n$ converges to $B$, and both of them converge absolutely, then their Cauchy product converges absolutely to $AB$.
\begin{proof}
PLACEHOLDER
\end{proof}
\sssc{Multiplication of power series}
If $f(x)=\sum_{n=0}^{\infty}a_nx^n$ and $g(x)=\sum_{n=0}^{\infty}b_nx^n$ for $x\in R$ then for $x\in R$,
\[f(x)g(x)=\sum_{n=0}^{\infty}a_nb_nx^n.\]
\bpr
PLACEHOLDER
\epr
\sssc{Division of power series}
If $f(x)=\sum_{n=0}^{\infty}a_nx^n$ and $g(x)=\sum_{n=0}^{\infty}b_nx^n$ for $x\in R$ and $b_0\neq 0$, then, with
\[c_0=\frac{a_0}{b_0},\]
\[c_n=\frac{1}{b_0}\qty(a_n-\sum_{k=1}^nb_kc_{n-k}),\]
for $x\in R$,
\[\frac{f(x)}{g(x)}=\sum_{n=0}^{\infty}c_nx^n.\]
\bpr
PLACEHOLDER
\epr
\ssc{p-test, p-series, $\frac{1}{x^p(\ln x)^q}$ test, and $\frac{1}{x^p(\ln x)^q}$ series}
\sssc{p-test}
For any fixed $a>0$, $\int_a^{\infty}\frac{1}{x^p}\dd{x}$ converges iff $p>1$, and if so,
\[\int_a^{\infty}\frac{1}{x^p}\dd{x}=\frac{1}{(p-1)a^{p-1}};\]
$\int_0^a\frac{1}{x^p}\dd{x}$ converges iff $p<1$, and if so,
\[\int_0^a\frac{1}{x^p}\dd{x}=-\frac{1}{(p-1)a^{p-1}}.\]
\begin{proof}
\[\int\frac{1}{x^p}\dd{x}=\begin{cases}
-\frac{1}{(p-1)x^{p-1}},\quad&p\neq 1\\
\ln(x),\quad&p=1
\end{cases}\]
\[\forall p<1\colon\nexists\lim_{x\to\infty}-\frac{1}{(p-1)x^{p-1}}\]
\[\nexists\lim_{x\to\infty}\ln(x)\]
\[\forall p>1\colon\lim_{x\to\infty}-\frac{1}{(p-1)x^{p-1}}=0\quad\int_a^{\infty}\frac{1}{x^p}\dd{x}=\frac{1}{(p-1)a^{p-1}}\]
\[\forall p>1\colon\nexists\lim_{x\to 0^+}-\frac{1}{(p-1)x^{p-1}}\]
\[\nexists\lim_{x\to 0^+}\ln(x)\]
\[\forall p<1\colon\lim_{x\to 0^+}-\frac{1}{(p-1)x^{p-1}}=0\quad\int_0^a\frac{1}{x^p}\dd{x}=-\frac{1}{(p-1)a^{p-1}}\]
\end{proof}
\sssc{p-series}
For any fixed $a\in\bbN$, $\sum_{n=a}^{\infty}\frac{1}{n^p}$, called p-series, converges iff $p>1$.
\bpr
For $p\le0$, it's obviously $\infty$. For $p>0$, $f(x)=\frac{1}{x^p}$ is monotone decreasing on $[1,\infty)$. By integral test and p-test, the statement holds.
\epr
\sssc{$\frac{1}{x^p(\ln x)^q}$ test}
For any fixed $a>1$, $\int_a^{\infty}\frac{1}{x^p(\ln x)^q}\dd{x}$ converges iff $p>1\lor(p=1\land q>1)$.

For any fixed $0<a<1$,
$\int_0^a\frac{1}{x^p(\ln x)^q}\dd{x}$ converges iff $p<1\lor(p=1\land q<1)$.
\begin{proof}
For $p>1$, take a $\varepsilon$ such that $p-1>\varepsilon>0$. $\frac{1}{x^{p-\varepsilon}}\ge\frac{1}{x^p(\ln x)^q}$ for all $x\ge b$ and $\frac{1}{x^{p-\varepsilon}}\le\frac{1}{x^p(\ln x)^q}$ for all $0<x\le b$ for some $b>0$. $\int_a^{\infty}\frac{1}{x^p(\ln x)^q}\dd{x}$ converges by p-test. Thus, by comparison theorem, $\int_a^{\infty}\frac{1}{x^p(\ln x)^q}\dd{x}$ converges. $\int_0^a\frac{1}{x^p(\ln x)^q}\dd{x}$ diverges by p-test. Thus, by comparison theorem, $\int_0^a\frac{1}{x^p(\ln x)^q}\dd{x}$ diverges.

For $p<1$, take a $\varepsilon$ such that $1-p>\varepsilon>0$. $\frac{1}{x^{p-\varepsilon}}\le\frac{1}{x^p(\ln x)^q}$ for all $x\ge b$ and $\frac{1}{x^{p-\varepsilon}}\ge\frac{1}{x^p(\ln x)^q}$ for all $0<x\le b$ for some $b>0$. $\int_a^{\infty}\frac{1}{x^p(\ln x)^q}\dd{x}$ diverges by p-test. Thus, by comparison theorem, $\int_a^{\infty}\frac{1}{x^p(\ln x)^q}\dd{x}$ diverges. $\int_0^a\frac{1}{x^p(\ln x)^q}\dd{x}$ converges by p-test. Thus, by comparison theorem, $\int_0^a\frac{1}{x^p(\ln x)^q}\dd{x}$ converges.

For $p=1$,
\[u=\ln x,\quad\dd{u}=\frac{1}{x}\dd{x}.\]
\[\int_a^{\infty}\frac{1}{x(\ln x)^q}\dd{x}=\int_{\ln a}^{\infty}\frac{1}{u^q}\dd{u},\]
which converges iff $q>1$ by p-test;
\[\int_0^a\frac{1}{x(\ln x)^q}\dd{x}=\int_0^{\ln a}\frac{1}{u^q}\dd{u},\]
which converges iff $q<1$ by p-test.
\end{proof}
\sssc{$\frac{1}{x^p(\ln x)^q}$ series}
For any fixed $a\in\bbN$, $\sum_{n=a}^{\infty}\frac{1}{n^p(\ln n)^q}$ converges iff $p>1\lor(p=1\land q>1)$.
\bpr
For $p\le0$, it's obviously $\infty$. For $p>0$, $f(x)=\frac{1}{x^p(\ln x)^q}$ is monotone decreasing on $[b,\infty)$ for some $b>1$. By integral test and $\frac{1}{x^p(\ln x)^q}$ test, the statement holds.
\epr
\ssc{Improper integral theorems}
\sssc{Convergence of improper integral of monotone increasing/decreasing integrand implies non-positive/non-negative integrand}
If function $f\colon[a,\infty)\to\bbR$ for some $a\in\bbR$ is monotone increasing and such that $\int_a^{\infty}f(x)\dd{x}$ converges, then $f(x)$ is non-positive.
\bpr
Use proof by contradiction. Suppose there exists $x_0>a$ such that $f(x_0)>0$. Then $f(x)>0$ for all $x\ge x_0$ and
\[\int_{x_0}^{\infty}f(x)\dd{x}=\infty.\]
Since $f(x)$ is monotone on $[a,x_0]$, it's BV and bounded on it, thus $\int_a^{x_0}f(x)\dd{x}\in\bbR$.
\epr
If function $f\colon[a,\infty)\to\bbR$ is monotone decreasing and such that $\int_a^{\infty}f(x)\dd{x}$ converges, then $f(x)$ is non-negative.
\sssc{Convergence of improper integral of monotone integrand implies limit of product of integrand and independent variable converges to 0}
If function $f\colon[a,\infty)\to\bbR$ for some $a\in\bbR$ is motone and $\int_a^{\infty}f(x)\dd{x}$ converges, then $\lim_{x\to\infty}xf(x)=0$. (And $\lim_{x\to\infty}f(x)=0$.)
\begin{proof}
$f(x)$ is either non-negative and non-increasing or non-positive and non-decreasing. Consider the former case.

Since $\int_a^{\infty}f(x)\dd{x}$ and $f(x)$ is non-negative, $\forall\varepsilon>0$, there exists $M>a$ such that $\forall A>M$:
\[0\leq\int_A^{\infty}f(x)\dd{x}<\varepsilon.\]
Thus for $t>2M$,
\[0\leq\int_{t/2}^tf(x)\dd{x}=\int_{t/2}^{\infty}f(x)\dd{x}-\int_t^{\infty}f(x)\dd{x}<2\varepsilon.\]
Since $f(x)$ is non-increasing,
\[f(x)\geq f(t),\quad\forall x\in[t/2,t]\]
\[\int_{t/2}^tf(x)\dd{x}\geq\int_{t/2}^tf(t)\dd{x}=\frac{t}{2}f(t)\]
\[tf(t)<4\varepsilon.\]
Thus, $\forall\varepsilon>0$, there exists $2M>0$ such that $\forall t>2M$,
\[tf(t)<\varepsilon,\]
that is,
\[\lim_{x\to\infty}xf(x)=0.\]
\end{proof}
\ssc{Absolute convergence implies convergence}
\sssc{For real series}
A real series is convergent if it is absolutely convergent.
\bpr
Suppose $\sum_{n=1}^{\infty}\abs{a_n}$ is convergent.
\[0\leq\sum_{n=1}^{\infty}a_n+\abs{a_n}\leq\sum_{n=1}^{\infty}2\abs{a_n}\]
By comparison test theorem, $\sum_{n=1}^{\infty}a_n+\abs{a_n}$ is convergent.
\[\sum_{n=1}^{\infty}a_n+\abs{a_n}-\sum_{n=1}^{\infty}\abs{a_n}=\sum_{n=1}^{\infty}a_n\]
\epr
\sssc{For improper integral}
An improper integral is convergent if it is absolutely convergent.
\bpr
Suppose
\[\int_a^b\abs{f(x)}\dd{x}\]
is a convergent improper integral and
\[\int_a^c\abs{f(x)}\dd{x}\]
is proper integral for any $a<c<b$.
\[0\leq\int_a^bf(x)+\abs{f(x)}\dd{x}\leq\int_a^b2\abs{f(x)}\dd{x}\]
By comparison test theorem, $\int_a^bf(x)+\abs{f(x)}\dd{x}$ is convergent.
\[\int_a^bf(x)+\abs{f(x)}\dd{x}-\int_a^b\abs{f(x)}\dd{x}=\int_a^bf(x)\dd{x}\]
\epr
\ssc{Deleting head}
\sssc{For real series}
For real sequence $\langle a_i\rangle_{i=m}^{\infty}$, $\sum_{i=m}^{\infty}a_i$ converges iff $\sum_{i=m+n}^{\infty}a_i$ converges, for all $n\in\bbN$.
\sssc{For improper integral}
If function $f\colon [a,b)\to\bbR$ is measurable with $b\in(a,\infty]$, $\int_a^bf(x)\dd{x}$ is improper integral, and $\int_a^cf(x)\dd{x}$ is proper integral for any $a<c<b$, then $\int_a^bf(x)\dd{x}$ converges iff $\int_c^bf(x)\dd{x} converges, for all $a<c<b$.
\ssc{Term test}
\sssc{For real series}
If $\sum_{n=0}^{\infty}a_n$ converges, then $\lim_{n\to\infty}a_n=0$.
\bpr
Let
\[S_n=\sum_{k=0}^na_k\]
and $S_{-1}=0$.
\[a_n=S_n-S_{n-1}.\]
Suppose $S_n$ converges to $S$. Then $S_{n-1}$ converges to $S$. Then $S_n-S_{n-1}$ converges to $S-S=0$.
\epr
\sssc{For improper integral}
If function $f\colon [a,\infty)\to\bbR$ is measurable and of bounded variation, $\int_a^bf(x)\dd{x}$ is proper integral for any $b>a$, and $\int_a^{\infty}f(x)\dd{x}$ converges to some real number, then
\[\lim_{x\to\infty}f(x)=0.\]
\bpr
Let $f^+$ and $f^-$ be the Jordan decomposition of $f$ such that $f=f^+-f^-$.
\[\int_a^{\infty}f(x)\dd{x}=\int_a^{\infty}f^+(x)\dd{x}-\int_a^{\infty}f^-(x)\dd{x}\]
and that LHS converges iff both terms of RHS converges.

$\lim_{x\to\infty}f(x)=0$ iff $\lim_{x\to\infty}f^+(x)=\lim_{x\to\infty}f^-(x)$.

Assume for contradiction that there exists $\varepsilon>0$ and $M\geq a$ such that
\[\forall x\geq M:\,f^+(x)>\varepsilon.\]
Then, $\int_M^{\infty}f^+(x)$ diverges.

Assume for contradiction that there exists $\varepsilon>0$ and $M\geq a$ such that
\[\forall x\geq M:\,f^-(x)>\varepsilon.\]
Then, $\int_M^{\infty}f^-(x)$ diverges.
\epr
\ssc{(Direct) comparison theorem}
\sssc{For real series}
If real sequences $\langle a_n\rangle_{n=1}^{\infty}$ and $\langle b_n\rangle_{n=1}^{\infty}$ are such that $a_n\geq b_n\geq 0$ for all $n\in\bbN$, then
\ben
\item $\sum_{n=1}^{\infty}a_n$ is convergent implies $\sum_{n=1}^{\infty}b_n$ is convergent.
\item $\sum_{n=1}^{\infty}b_n$ is divergent implies $\sum_{n=1}^{\infty}a_n$ is divergent.
\een
\begin{proof}
Let $\left\langle s_n=\sum_{i=1}^na_i\right\rangle_{i=1}^{\infty}$ and $\left\langle t_n=\sum_{i=1}^nb_i\right\rangle_{i=1}^{\infty}$ be the sequences of partial sums of $\langle a_n\rangle_{n=1}^{\infty}$ and $\langle b_n\rangle_{n=1}^{\infty}$.

$s_n$ and $t_n$ are monotone increasing for all $n\in\bbN$ since $a_n,b_n\geq 0$. $s_n\geq t_n$ for all $n\in\bbN$ since $a_n\geq b_n$.

Statement 1:
\[\sum_{n=1}^{\infty}a_n=\lim_{n\to\infty}s_n=L.\]
Since
\[0\leq t_n\leq s_n,\]
then,
\[\sum_{n=1}^{\infty}b_n=\lim_{n\to\infty}t_n\]
exists and $\in[0,L]$.

Statement 2 is the contrapositive of statement 1.
\end{proof}
\sssc{For improper integral}
If functions $f\colon [a,b)\to\bbR$ and $g\colon [a,b)\to\bbR$ are measurable with $b\in(a,\infty]$, $\int_a^bf(x)\dd{x}$ and $\int_a^bg(x)\dd{x}$ are improper integrals, $\int_a^cf(x)\dd{x}$ and $\int_a^cg(x)\dd{x}$ are proper integrals for any $a<c<b$,
and $f(x)\geq g(x)\geq 0$ a.e. for $x\in[a,b)$, then
\ben
\item $\int_a^bf(x)\dd{x}$ is convergent implies $\int_a^bg(x)\dd{x}$ is convergent.
\item $\int_a^bg(x)\dd{x}$ is divergent implies $\int_a^bf(x)\dd{x}$ is divergent.
\een
\begin{proof}
Let
\[F(t)=\int_a^tf(x)\dd{x},\quad t\in[a,b),\]
\[G(t)=\int_a^tg(x)\dd{x},\quad t\in[a,b).\]
$F$ and $G$ are monotone increasing since $f(x),g(x)\geq 0$ a.e. for $x\in[a,b)$.

$F(x)\geq G(x)$ since $f(x)\geq g(x)$ a.e. for $x\in[a,b)$.

Statement 1:
\[\int_a^bf(x)\dd{x}=\lim_{t\to b}F(t)=L.\]
Since
\[0\leq G(x)\leq F(x),\]
then,
\[\int_a^bg(x)\dd{x}=\lim_{t\to b}G(t)\]
exists and $\in [0,L]$.

Statement 2 is the contrapositive of statement 1.
\end{proof}
\ssc{Limit comparison theorem (LCT)}
\sssc{For real series}
If real sequences $\langle a_n\rangle_{n=1}^{\infty}$ and $\langle b_n\rangle_{n=1}^{\infty}$ are such that $a_n\geq 0$ and $b_n>0$ for all $n\geq m$ for some $m\in\bbN$, if
\[L=\lim_{n\to\infty}\frac{a_n}{b_n}\in[0,\infty],\]
then:
\ben
\item If $0<L<\infty$, then $\sum_{n=1}^{\infty}a_n$ is convergent iff $\sum_{n=1}^{\infty}b_n$ is convergent. (classical LCT)
\item if $L=0$, then:
\ben
\item $\sum_{n=1}^{\infty}b_n$ is convergent implies $\sum_{n=1}^{\infty}a_n$ is convergent.
\item $\sum_{n=1}^{\infty}a_n$ is divergent implies $\sum_{n=1}^{\infty}b_n$ is divergent.
\een
(one-sided LCT)
\item If $L=\infty$, then:
\ben
\item $\sum_{n=1}^{\infty}a_n$ is convergent implies $\sum_{n=1}^{\infty}b_n$ is convergent.
\item $\sum_{n=1}^{\infty}b_n$ is divergent implies $\sum_{n=1}^{\infty}a_n$ is divergent.
\een
(one-sided LCT)
\een
\begin{proof}
Statement 1:

For every $\varepsilon>0$, there exists $N\in\bbN$ such that for all $n\geq N$, $\abs{\frac{a_n}{b_n}-L}<\varepsilon$, equivalently,
\[(L-\varepsilon)b_n<a_n<(L+\varepsilon)b_n.\]
As $L>0$, choose $\varepsilon$ such that $L-\varepsilon>0$, so that
\[\frac{a_n}{L+\varepsilon}<b_n<\frac{a_n}{L-\varepsilon}.\]
$\sum_{n=1}^{\infty}\frac{a_n}{L+\varepsilon}$ is convergent iff $\sum_{n=1}^{\infty}a_n$ is convergent iff $\sum_{n=1}^{\infty}\frac{a_n}{L-\varepsilon}$ is convergent.

And by direct comparison test,
\bit
\item $\sum_{n=1}^{\infty}\frac{a_n}{L-\varepsilon}$ is convergent implies $\sum_{n=1}^{\infty}b_n$ is convergent,
\item $\sum_{n=1}^{\infty}b_n$ is convergent implies $\sum_{n=1}^{\infty}\frac{a_n}{L+\varepsilon}$ is convergent.
\eit

Statement 2:

There exists $N\in\bbN$ such that for all $n\geq N$, $0\leq\frac{a_n}{b_n}<1$, equivalently,
\[0\leq a_n<b_n.\]

By direct comparison test of $\sum_{n=N}^{\infty}a_n$ and $\sum_{n=N}^{\infty}b_n$, the statements hold.

Statement 3:

There exists $N\in\bbN$ such that for all $n\geq N$, $\frac{a_n}{b_n}>1$, equivalently,
\[a_n>b_n>0.\]

By direct comparison test of $\sum_{n=N}^{\infty}a_n$ and $\sum_{n=N}^{\infty}b_n$, the statements hold.
\end{proof}
\sssc{For improper integral}
If functions $f\colon [a,b)\to\bbR$ and $g\colon [a,b)\to\bbR$ are measurable with $b\in(a,\infty]$, $\int_a^bf(x)\dd{x}$ and $\int_a^bg(x)\dd{x}$ are improper integrals, $\int_a^cf(x)\dd{x}$ and $\int_a^cg(x)\dd{x}$ are proper integrals for any $a<c<b$, $f(x)\geq 0$ and $g(x)>0$ for $x\in[a,b)$, if
\[L=\lim_{x\to b^-}\frac{f(x)}{g(x)}\in[0,\infty],\]
then:
\ben
\item If $0<L<\infty$, then $\int_a^bf(x)\dd{x}$ is convergent iff $\int_a^bg(x)\dd{x}$ is convergent. (classical LCT)
\item if $L=0$, then:
\ben
\item $\int_a^bg(x)\dd{x}$ is convergent implies $\int_a^bf(x)\dd{x}$ is convergent.
\item $\int_a^bf(x)\dd{x}$ is divergent implies $\int_a^bg(x)\dd{x}$ is divergent.
\een
(one-sided LCT)
\item If $L=\infty$, then:
\ben
\item $\int_a^bg(x)\dd{x}$ is divergent implies $\int_a^bf(x)\dd{x}$ is divergent.
\item $\int_a^bf(x)\dd{x}$ is convergent implies $\int_a^bg(x)\dd{x}$ is convergent.
\een
(one-sided LCT)
\een
\begin{proof}
Statement 1:

For every $\varepsilon>0$, there exists $a\leq M<b$ such that for all $b>x\geq M$, $\abs{\frac{f(x)}{g(x)}-L}<\varepsilon$, equivalently,
\[(L-\varepsilon)g(x)<f(x)<(L+\varepsilon)g(x).\]
As $L>0$, we can choose $\varepsilon$ such that $L-\varepsilon>0$, so that
\[\frac{f(x)}{L+\varepsilon}<g(x)<\frac{f(x)}{L-\varepsilon}.\]
$\int_a^b\frac{f(x)}{L+\varepsilon}\dd{x}$ is convergent iff $\int_a^bf(x)\dd{x}$ is convergent iff $\int_a^b\frac{f(x)}{L-\varepsilon}\dd{x}$ is convergent.

And by direct comparison test,
\bit
\item $\int_a^b\frac{f(x)}{L-\varepsilon}\dd{x}$ is convergent implies $\int_a^bg(x)\dd{x}$ is convergent,
\item $\int_a^bg(x)\dd{x}$ is convergent implies $\int_a^b\frac{f(x)}{L+\varepsilon}\dd{x}$ is convergent.
\eit

Statement 2:

There exists $a\leq M<b$ such that for all $b>x\geq M$, $0\leq\frac{f(x)}{g}<1$, equivalently,
\[0\leq f(x)<g(x).\]

By direct comparison test of $\int_M^bf(x)\dd{x}$ and $\int_M^bg(x)\dd{x}$, the statements hold.

Statement 3:

There exists $a\leq M<b$ such that for all $b>x\geq M$, $\frac{f(x)}{g(x)}>1$, equivalently,
\[f(x)>g(x)>0.\]

By direct comparison test of $\int_M^bf(x)\dd{x}$ and $\int_M^bg(x)\dd{x}$, the statements hold.
\end{proof}
\ssc{Dirichlet test or Abel's test}
\sssc{For real series}
If $\langle a_n\rangle_{n=1}^{\infty}$ is a sequence such that the sequence of partial sums $\langle s_n=\bcs0,\quad&n=0\\\sum_{i=0}^na_i,\quad&n\in\bbN\ecs\rangle_{n=0}^{\infty}$ is bounded and $\langle b_n\rangle_{n=1}^{\infty}$ is a monotone sequence such that $\lim_{n\to\infty}b_n=0$, then
\[\sum_{n=1}^{\infty}a_nb_n\]
converges.
\begin{proof}
Choose $M$ such that $\langle\abs{s_n}\rangle_{n=0}^{\infty}$ is bounded above by $M$.

For any $\varepsilon>0$, there is $N\in\bbN$ such that $\abs{b_m}<\frac{\varepsilon}{2M}$ for all $m\geq N$.

For any $N\leq p\leq q$, using Abel transformation,
\[\ba
\abs{\sum_{n=p}^qa_nb_n}&=\abs{\sum_{n=p}^q\qty(s_{n+1}-s_n)b_n}\\
&=\abs{s_{q+1}b_{q+1}-s_pb_p-\sum_{n=p}^qs_{n+1}\qty(b_{n+1}-b_n)}\\
&\leq M\qty(\abs{b_{q+1}}+\abs{b_p}+\sum_{n=p}^q\abs{b_{n+1}-b_n})\\
&=M\qty(\abs{b_{q+1}}+\abs{b_p}+\abs{b_{q+1}-b_p})\\
&=M\qty(\abs{b_{q+1}}+\abs{b_p}+\abs{b_p}-\abs{b_{q+1}})\\
&=2M\abs{b_p}\\
&\leq\varepsilon
\ea\]
\end{proof}
\sssc{For improper integral}
If $f\colon [a,b)\to\bbR$ with $b\in(a,\infty]$ is a function such that the function
\[F(x)=\int_a^xf(t)\dd{t},\quad b>x\geq a\]
is bounded, $g\colon [a,b)\to\bbR$ is a monotone and measurable function such that $\lim_{x\to b^-}g(x)=0$, then
\[\int_a^bf(x)g(x)\dd{x}\]
converges.
\begin{proof}
Rewrite as Riemann–Stieltjes integral
\[\int_a^bf(x)g(x)\dd{x}=\int_{x=a}^bg(x)\dd{F(x)}.\]

Choose $M$ such that $\abs{F(x)}$ is bounded above by $M$.

For any $\varepsilon>0$, there is $a\leq N<b$ such that $\abs{g(N)}<\frac{\varepsilon}{2M}$ for all $m\geq N$.

For any $N\leq p\leq q<b$, using integration by parts of Riemann–Stieltjes integral,
\[\ba
\abs{\int_{x=p}^qg(x)\dd{F(x)}}&=\abs{F(q)g(q)-F(p)g(p)-\int_{x=p}^qf(x)\dd{g(x)}}\\
&\leq M\qty(\abs{g(q)}+\abs{g(p)}+\int_{x=p}^q\dd{g(x)})\\
&=M\qty(\abs{g(q)}+\abs{g(p)}+\abs{g(q)-g(p)})\\
&=M\qty(\abs{g(q)}+\abs{g(p)}+\abs{g(p)}-\abs{g(q)})\\
&=2M\abs{g(p)}\\
&\leq\varepsilon
\ea\]
\end{proof}
\ssc{Cesàro-summability theorem}
\sssc{For real sequence}
If a real sequence converges to $L\in\bbR$, then the sequence is Cesàro-summable and its Cesàro sum is $L$.
\begin{proof}
PLACEHOLDER
\end{proof}
\sssc{For real function}
If a real function $f$ and $a\in\bbR$ is such that $\lim_{x\to\infty}f(x)=L\in\bbR$ and $\int_a^bf(x)\dd{x}$ is a proper integral for all $b\in[a,\infty)$, then the function is Cesàro-summable over $[a,\infty)$ and its Cesàro sum over $[a,\infty)$ is $L$.
\begin{proof}
PLACEHOLDER
\end{proof}
If a real function $f$ and $b\in\bbR$ is such that $\lim_{x\to-\infty}f(x)=L\in\bbR$ and $\int_a^bf(x)\dd{x}$ is a proper integral for all $a\in(-\infty,b)$, then the function is Cesàro-summable over $(-\infty,b]$ and its Cesàro sum over $(-\infty,b]$ is $L$.